\heading{数学物理方法（下）-模拟题1-期末}

\noteinfo{题目和答案由AI生成。}

\begin{question}
计算卷积积分
\[
I(x)=\int_0^x \sin(x-t)J_0(t)\,dt,
\]
其中 $J_0$ 为第一类零阶 Bessel 函数。
\begin{solution}

用 Laplace 变换。已知
\[
\mathcal L\{\sin x\}=\frac{1}{s^2+1},\qquad \mathcal L\{J_0(x)\}=\frac{1}{\sqrt{s^2+1}}.
\]
故
\[
\mathcal L\{I\}=\frac{1}{(s^2+1)^{3/2}}.
\]
又利用恒等式
\[
\frac{d}{dx}\bigl(xJ_1(x)\bigr)=xJ_0(x)
\]
可知
\[
\mathcal L\{1-J_0(x)\}=\frac{1}{s}-\frac{1}{\sqrt{s^2+1}}
=\frac{\sqrt{s^2+1}-s}{s\sqrt{s^2+1}},
\]
不够直接；更直接地注意到
\[
\mathcal L\left\{\int_0^x tJ_0(t)\,dt\right\}=\frac1s\mathcal L\{xJ_0(x)\}=\frac1s\frac{s}{(s^2+1)^{3/2}}=\frac{1}{(s^2+1)^{3/2}}.
\]
因此
\[
I(x)=\int_0^x tJ_0(t)\,dt=xJ_1(x).
\]
故
\ansmath{\int_0^x \sin(x-t)J_0(t)\,dt=xJ_1(x).}
\end{solution}
\end{question}

\begin{question}
在圆盘 $x^2+y^2<a^2$ 中求解泊松方程
\[
-\kappa\nabla^2u=y,
\qquad u|_{x^2+y^2=a^2}=x^2-y^2.
\]
要求写出一个显式特解，并说明一般解结构。
\begin{solution}

先取边界项的调和延拓
\[
 u_1=x^2-y^2.
\]
再找零边界特解。注意到二维中
\[
\nabla^2(r^2y)=8y,
\qquad r^2=x^2+y^2.
\]
因此可取
\[
 u_2=-\frac{(r^2-a^2)y}{8\kappa}.
\]
因为在边界 $r=a$ 时 $u_2=0$，且
\[
-\kappa\nabla^2u_2=y.
\]
故一个显式解为
\ansmath{u(x,y)=x^2-y^2-\frac{(x^2+y^2-a^2)y}{8\kappa}.}
若只要求满足同样边界数据，则这已是唯一解；若题目改成带时间项的热方程，则一般解应写成“该稳态解 $+$ 圆盘上零边界齐次热方程的 Bessel 本征展开”。
\end{solution}
\end{question}

\begin{question}
给定算符
\[
\hat L=-\frac{d}{dx},
\qquad D(\hat L)=\{u\in L^2(0,\pi):u(\pi)=\gamma u(0)\},
\]
求伴算符 $\hat L^\dagger$，并求何时有 $\hat L^\dagger=-\hat L$。在该条件下求全部本征值和本征函数。
\begin{solution}

对 $u\in D(\hat L)$、$v\in D(\hat L^\dagger)$，
\[
\langle \hat Lu,v\rangle=-\int_0^{\pi}u'\bar v\,dx=
\int_0^{\pi}u\overline{v'}\,dx-u(\pi)\overline{v(\pi)}+u(0)\overline{v(0)}.
\]
利用 $u(\pi)=\gamma u(0)$，边界项为
\[
u(0)\bigl(\overline{v(0)}-\gamma\overline{v(\pi)}\bigr).
\]
对任意 $u$ 消失，当且仅当
\[
 v(0)=\overline\gamma\,v(\pi).
\]
故
\ansmath{\hat L^\dagger=\frac{d}{dx},\qquad D(\hat L^\dagger)=\{v:v(0)=\overline\gamma v(\pi)\}.}
要使 $\hat L^\dagger=-\hat L$，只需定义域相同，即
\[
 v(\pi)=\gamma v(0)\iff v(0)=\overline\gamma v(\pi),
\]
这等价于
\ansmath{|\gamma|=1.}
设 $\gamma=e^{i\theta}$。本征方程
\[
 v'=\lambda v
\]
给出
\[
 v(x)=Ce^{\lambda x}.
\]
代入 $v(0)=e^{-i\theta}v(\pi)$，得
\[
1=e^{-i\theta}e^{\lambda\pi}\Longrightarrow e^{\lambda\pi}=e^{i\theta}.
\]
故
\ansmath{\lambda_n=\frac{i(\theta+2\pi n)}{\pi},\qquad v_n(x)=e^{\lambda_n x},\qquad n\in\mathbb Z.}
\end{solution}
\end{question}

\begin{question}
求解无界域热方程
\[
\begin{cases}
 u_t-\kappa u_{xx}=e^{-\alpha t}\delta(x), & -\infty<x<\infty,\ t>0,\\
 u(x,0)=0.
\end{cases}
\]

\begin{solution}

利用热核
\[
G(x,t)=\frac{1}{\sqrt{4\pi\kappa t}}\exp\!\left(-\frac{x^2}{4\kappa t}\right),\qquad t>0.
\]
Duhamel 原理给出
\[
 u(x,t)=\int_0^t e^{-\alpha\tau}G(x,t-\tau)\,d\tau.
\]
因此
\ansmath{u(x,t)=\int_0^t \frac{e^{-\alpha\tau}}{\sqrt{4\pi\kappa (t-\tau)}}
\exp\!\left(-\frac{x^2}{4\kappa (t-\tau)}\right)d\tau.}
若换元 $s=t-\tau$，也可写成
\[
 u(x,t)=e^{-\alpha t}\int_0^t \frac{e^{\alpha s}}{\sqrt{4\pi\kappa s}}
\exp\!\left(-\frac{x^2}{4\kappa s}\right)ds.
\]
这就是所求显式积分表示。
\end{solution}
\end{question}

\begin{question}
把半径为 $a$、总电量为 $Q$ 的均匀带电圆环放在平面 $z=0$ 上、圆心在原点。分别写出其在直角坐标、柱坐标、球坐标下的电荷密度分布。
\begin{solution}

圆环是线电荷，线密度
\[
\lambda=\frac{Q}{2\pi a}.
\]
\begin{enumerate}[label=（\arabic*）,leftmargin=2.8em,itemsep=0.2em]
\item 在柱坐标下最简洁：圆环满足 $\rho=a,\ z=0$，故
\ansmath{\rho_e(\rho,\varphi,z)=\frac{Q}{2\pi a}\,\delta(\rho-a)\delta(z).}
注意体元为 $dV=\rho\,d\rho\,d\varphi\,dz$，积分后恰给出 $Q$。
\item 在直角坐标下可写成
\ansmath{\rho_e(x,y,z)=\frac{Q}{2\pi a}\,\delta\!\left(\sqrt{x^2+y^2}-a\right)\delta(z).}
也可用 $\delta(x^2+y^2-a^2)$ 改写，但需同时补上雅可比因子。
\item 在球坐标中圆环对应 $r=a,\ \theta=\pi/2$。因体元为 $r^2\sin\theta\,dr\,d\theta\,d\varphi$，可取
\ansmath{\rho_e(r,\theta,\varphi)=\frac{Q}{2\pi a^2}\,\delta(r-a)\,\delta\!\left(\theta-\frac\pi2\right).}
\end{enumerate}
\end{solution}
\end{question}

\begin{question}
在半圆域
\[
0<r<R,\qquad 0<\theta<\pi
\]
上写出满足零边界条件
\[
G|_{\theta=0}=G|_{\theta=\pi}=G|_{r=R}=0
\]
的 Laplace 算符 Green 函数的本征展开结构。提示：先做分离变量。
\begin{solution}

零边界条件下，角向本征函数为
\[
\Theta_n(\theta)=\sin n\theta,\qquad n=1,2,\dots
\]
径向部分满足 Bessel 方程，且在 $r=0$ 有界、在 $r=R$ 为零，因此应取
\[
R_{n,m}(r)=J_n\!\left(\frac{j_{n,m}}{R}r\right),
\qquad J_n(j_{n,m})=0,
\qquad m=1,2,\dots
\]
于是 Green 函数可按完整正交系展开为
\ansmath{G(r,\theta;r',\theta')=
\sum_{n=1}^{\infty}\sum_{m=1}^{\infty}
\frac{C_{n,m}}{\lambda_{n,m}}
J_n\!\left(\frac{j_{n,m}}{R}r_<\right)
J_n\!\left(\frac{j_{n,m}}{R}r_>\right)
\sin n\theta\sin n\theta',}
其中
\[
\lambda_{n,m}=\left(\frac{j_{n,m}}{R}\right)^2,
\]
$C_{n,m}$ 由正交归一化决定。更标准地写成
\[
G=\sum_{n,m}
\frac{J_n(j_{n,m}r/R)J_n(j_{n,m}r'/R)\sin n\theta\sin n\theta'}{
\lambda_{n,m}\,N_{n,m}},
\]
其中
\[
N_{n,m}=\int_0^{\pi}\!\sin^2 n\theta\,d\theta\int_0^R rJ_n^2(j_{n,m}r/R)\,dr.
\]
这就是所求本征展开结构。
\end{solution}
\end{question}


